\section{Why it might work --- and what we could not show}
\label{sec:mechanism}

This section is deliberately weaker than the two before it, and it reports a negative result
about our own preferred explanation. Three candidate mechanisms; the full analysis is in
Appendix~\ref{app:mechanism}.

\textbf{A more faithful reward? Over the run, yes; at a matched policy, no.} Measuring pairwise
rank agreement between the partial-conversation training reward and the full-session score, the
$\Kf$ signal is the more faithful ranker pooled over the run ($0.909$ vs.\ $0.873$ under the
training oracle, $0.800$ vs.\ $0.747$ held out) --- but the iteration-level Wilcoxon does
\emph{not} clear $.05$, so the difference is small-and-consistent, not significant, and over the
run the two arms are different policies on different conversations, so this cut cannot establish a
mechanism. The one cut where the scoring rule is close to the only difference (iteration 1, both
arms branching from the same base policy) shows \textbf{nothing}: all four deltas straddle zero,
and the per-length-bin breakdown points in opposite directions under the two graders --- the shape
of no effect. The proxy-quality mechanism is \textbf{unsupported}.

\textbf{Sharper group discrimination? No --- a pure rescaling.} Look-ahead spreads the eight
siblings' rewards further apart (best-minus-worst margin $1.300\times$), but the within-group
standard deviation grows by almost the same factor ($1.293\times$), and GRPO divides the advantage
by it. The ratio of ratios is $1.006$ $[1.002, 1.010]$ --- a rescaling the advantage normalisation
absorbs. Look-ahead changes \emph{which} candidate the group prefers, not how confidently.

\textbf{A different update direction? Barely.} Writing each iteration's update direction as the
advantage-weighted sum of its candidates' sentence embeddings, the two arms' pooled directions
have a cosine of $0.804$ ($0.851$ attenuation-corrected). The $\Kz$ and $\Kf$ policies are pushed
along substantially the \emph{same} semantic direction.

That is the puzzle this paper leaves open, stated sharply: the same optimizer, pushed in nearly the
same direction, with a reward that is not measurably better conditioned and not measurably more
faithful at a fixed policy, nevertheless ends more than twice as far from base. What differs is
\emph{which} of eight siblings is preferred at each step --- a small, repeated difference in
selection, compounded over ten iterations of a loop that regenerates its own training data from
the policy it just updated. Isolating that pathway needs a cut this experiment does not contain:
the same policy trained under both rewards with the horizon swapped mid-run.
